\section{Introduction}
\label{sec:introduction}

Large language models (LLMs) are increasingly deployed as components of
multi-agent systems in which models critique, debate, verify, or revise one
another's outputs. These systems can improve factuality and reasoning by
combining multiple model responses \cite{du2024debate}. Their success,
however, depends on more than whether a Critic can produce useful feedback. A
Solver must also decide whether that feedback should change its answer. Blind
deference can propagate incorrect recommendations, while indiscriminate
resistance can prevent useful correction. Prior work on self-correction and
sycophancy shows that externally prompted revision does not necessarily imply
reliable error correction \cite{huang2024selfcorrect,sharma2024sycophancy}.

Behavioral evaluation can reveal whether a Solver changed its answer, but it
cannot by itself identify the internal computations responsible for that
decision. A representation may correlate with feedback adoption because it
encodes disagreement, answer content, Critic quality, or a downstream
consequence of the decision. Establishing a mechanism therefore requires two
steps: finding internal features that predict behavior on held-out data and
intervening on those features to test whether they influence behavior.

Sparse autoencoders (SAEs) offer a practical representation for this analysis.
They decompose dense neural activations into a larger set of sparsely active
features, some of which can be more interpretable than individual model
dimensions \cite{cunningham2023sparse,gao2025scaling}. We use SAEs to encode
the Solver's hidden states while it processes Critic feedback. Regularized
linear probes then test whether the resulting feature activations predict
whether the Solver adopts the Critic's recommendation or retains its original
answer. Finally, selected features are ablated or modified during a replayed
forward pass, allowing us to compare targeted interventions with random-feature
controls.

Our experiments use MuSiQue, a benchmark designed to require composition over
multiple reasoning steps \cite{trivedi2022musique}. For each question, a
Gemma 3 4B instruction-tuned Solver produces an initial answer and a Qwen3 4B
instruction-tuned Critic independently produces a review and recommendation.
The Solver is subsequently shown its prior answer and the Critic's candidate
answer in randomized order. This candidate-presentation protocol makes
adoption and retention more directly observable than unconstrained free-form
revision while preserving agreement and unresolved cases as separate outcomes.

This project addresses three primary research questions:

\begin{enumerate}
    \item Which internal Solver representations predict adoption versus
    retention of Critic feedback?
    \item Do interaction-trained SAEs isolate feedback-related structure more
    effectively than a general pretrained SAE?
    \item Do targeted interventions on probe-selected SAE features causally
    change feedback-uptake behavior relative to appropriate controls?
\end{enumerate}

The main contribution is an end-to-end, reproducible workflow connecting
behavioral data collection, leakage-aware partitioning, sparse representation
learning, predictive probing, and causal intervention. The workflow also
retains agreement, refusal, ambiguous, and malformed responses rather than
forcing every interaction into a binary target. This distinction is important:
a model cannot meaningfully adopt a Critic's answer when both agents already
agree, and uncertain text matches should not be treated as ground-truth
behavior labels.
